\documentclass[14pt,usenames,dvipsnames,svgnames]{beamer}
\usepackage[english]{babel}
\usepackage[misc]{ifsym}
\usepackage{mathspec,mathabx,multicol,listings,comment,fontawesome,pgfornament,comment,pifont}
\usepackage[compat=0.2]{TeXnicalities}
\usetikzlibrary{shapes,matrix}

\setsansfont{Yanone Kaffeesatz}[
    UprightFont     = *-Regular ,
    BoldFont        = *-Bold ,
    BoldItalicFont  = *-Bold ,
    BoldSlantedFont = *-Bold ,
    ItalicFont      = *-Light ,
    SlantedFont     = *-Light ,
    SmallCapsFont   = *-Thin
]

\graphicspath{{Figures/},{Figures/Emoji/}}
\usepackage{listings}
\def\transpPerc{100}
%listings set
\lstdefinestyle{MyCpp}{
% backgroundcolor=\color{white},    % choose the background color; you must add \usepackage{color} or \usepackage{xcolor}
breakatwhitespace=false,            % sets if automatic breaks should only happen at whitespace
breaklines=true,                    % sets automatic line breaking
captionpos=b,                       % sets the caption-position to bottom
deletekeywords={...},               % if you want to delete keywords from the given language
escapeinside={@|}{|@},                % if you want to add LaTeX within your code
extendedchars=true,                 % lets you use non-ASCII characters; for 8-bits encodings only,
                                    % does not work with UTF-8
frame=none  ,                       % adds a frame around the code
numbers=none,                       % where to put the line-numbers; possible values are (none, left, right)
numbersep=5pt,                      % how far the line-numbers are from the code
numberstyle=\tiny\color{black},     % the style that is used for the line-numbers
rulecolor=\color{black},            % if not set, the frame-color may be changed on line-breaks within not-black text
                                    % (e.g. comments (green here))
showspaces=false,                   % show spaces everywhere adding particular underscores; it overrides 'showstringspaces'
showstringspaces=false,             % underline spaces within strings only
showtabs=false,                     % show tabs within strings adding particular underscores
stepnumber=2,                       % the step between two line-numbers. If it's 1, each line will be numbered
stringstyle=\color{OliveGreen},     % string literal style
tabsize=2,                          % sets default tabsize to 2 spaces
title=\lstname,                     % show the filename of files included with \lstinputlisting; also try caption instead of title
%
%Base style for this presentation 
keepspaces=true,                    % keeps spaces in text, useful for keeping indentation of code
                                    % (possibly needs columns=flexible)
keywordstyle=\color{Cyan},          % keyword style
language=C++,
basicstyle=\ttfamily\scriptsize\color{black},
keywordstyle=\color{OliveGreen},
stringstyle=\color{Magenta},
commentstyle=\color{red},
moredelim=[is][\color{ForestGreen}]{|+}{+|},
literate=% literate={<replace>}{<replacement text>}{<width>}
  {\#define}{{{\color{CarnationPink}\#define}}}{6}
  {\#include}{{{\color{CarnationPink}\#include}}}{7},
morekeywords={real, vector, Minesweeper, myVecVec, Cell, SetOfCells, regex, T, string, CD,
              invalid_argument, Jukebox, Coffee, Version, Point2D, Point3D, Person, pair},
emph=[1]{calculateAverage, size, exit, begin, end, find, makeReportForUser, isElementInVector,
         main, WelcomeUserToTheGame, playGame, PrintGameResult, isPrime, sqrt, push_back, FC,
         getFlaggedCells, isFlagged, at, isEligibleForFullBenefits, gsl_fcmp, getDefaultForHelper,
         trim_right, addCD, MakeCoffee, WarmUpMachineIfNeeded, GrindCoffee, SetPressureAndTemperature,
         BrewCoffee, getMajorVersionNumber, getMinorVersionNumber, getBuildNumber, Distance, factorial,
         Q_rsqrt, Q_fastInverseSqrt, calculateApproximateInverseSqrt, refineResult, get, fetch, retrieve, obtain},
emphstyle=[1]{\color{NavyBlue}}, %Functions
emph=[2]{cout, data, sum, sum1, sum2, data, dataSet, dataSet1, dataSet2, first, last, argc, argv, minesweeper,
         number, i, c, f, b, list, gameBoard, flaggedCells, cell, realDaysPerIdealWeek, totalRealWeeksNeeded,
         realDaysPerTask, taskDaysEstimate, realWeeksPerTask, j, e, employee, timePattern, iterationNumber,
         functionValue, epsilon, it, value, keyToDecipher, longTitle, longAuthor, durationInMinutes, cd,
         typeOfCoffee, x, y, z, x2, approximateResult, floatAsInteger, initialNumber, refinedResult,
         threehalfs, result, timestamp, ymdhms},
emphstyle=[2]{\color{Orange}}, %Variables
emph=[3]{if, else, while, do, for, case, switch},
emphstyle=[3]{\color{violet}}, %Loops, if, etc.
emph=[4]{const, auto, break, continue, default, static, return, struct, NULL, sizeof, typedef,
         template, typename, true, false, throw, public, class},
emphstyle=[4]{\color{ProcessBlue}}, %Logical keywords
emph=[5]{std, boost},
emphstyle=[5]{\color{Maroon}}, %Namespaces
emph=[6]{STATUS_VALUE, FLAGGED, WORKING_DAYS_PER_WEEK, NUMBER_OF_TASKS, HOURLY_FLAG,
         CG_CHECK_RESIDUUM_EVERY, BOOST_AUTO_TEST_CASE, BOOST_TEST_MODULE, BOOST_REQUIRE_EQUAL},
emphstyle=[6]{\color{Gray}}, %Macros
emph=[7]{age, flags, title, author, tracks, duration, cdList},
emphstyle=[7]{\color{Peach!50!Purple}}, %Members
emph=[8]{Factorial, SpecialCases, NormalCases},
emphstyle=[8]{\color{Blue}}, %Boost cases and suite
}


\def\CPP{\lstinline[style=MyCpp, basicstyle=\ttfamily\color{black}]}

\makeatletter
\newenvironment{CenteredBox}{% 
\begin{Sbox}}{% Save the content in a box
\end{Sbox}\centerline{\parbox{\wd\@Sbox}{\TheSbox}}}% And output it centered
\makeatother


\usetikzlibrary{fadings}

\pgfkeys{%
/piechartthreed/.cd,
scale/.code                =  {\def\piechartthreedscale{#1}},
mix color/.code            =  {\def\piechartthreedmixcolor{#1}},
background color/.code     =  {\def\piechartthreedbackcolor{#1}},
name/.code                 =  {\def\piechartthreedname{#1}}}

\newcommand\piechartthreed[2][]{% 
   \pgfkeys{/piechartthreed/.cd,
     scale            = 1,
     mix color        = white,
     background color = white,
     name             = pc} 
  \pgfqkeys{/piechartthreed}{#1}
  \begin{scope}[scale=\piechartthreedscale] 
  \begin{scope}[xscale=5,yscale=3] 
     \path[preaction={fill=black,opacity=.8,
         path fading=circle with fuzzy edge 20 percent,
         transform canvas={yshift=-15mm*\piechartthreedscale}}] (0,0) circle (1cm);
    \fill[gray](0,0) circle (0.5cm);  
     \path[preaction={fill=\piechartthreedbackcolor,opacity=.8,
          path fading=circle with fuzzy edge 20 percent,
          transform canvas={yshift=-10mm*\piechartthreedscale}}] (0,0) circle (0.5cm);
     \pgfmathsetmacro\totan{0} 
     \global\let\totan\totan 
     \pgfmathsetmacro\bottoman{180} \global\let\bottoman\bottoman 
     \pgfmathsetmacro\toptoman{0}   \global\let\toptoman\toptoman 
     \begin{scope}[draw=black,thin]
     \foreach \an/\col [count=\xi] in {#2}{%
     \def\space{ } 
        \coordinate (\piechartthreedname\space\xi) at (\totan+\an/2:0.75cm); 
        \ifdim 180pt>\totan pt 
         \ifdim 0pt=\toptoman pt
            \shadedraw[left color=\col!50!\piechartthreedmixcolor,
                       right color=\col!25!\piechartthreedmixcolor,
                       draw=black,very thin] (0:.5cm) -- ++(0,-3mm) arc (0:\totan+\an:.5cm) 
                                                       -- ++(0,3mm)  arc (\totan+\an:0:.5cm);
            \pgfmathsetmacro\toptoman{180} 
            \global\let\toptoman\toptoman         
            \else
            \shadedraw[left color=\col!50!\piechartthreedmixcolor,
                       right color=\col!25!\piechartthreedmixcolor,
                       draw=black,very thin](\totan:.5cm)-- ++(0,-3mm) arc(\totan:\totan+\an:.5cm)
                                                        -- ++(0,3mm)  arc(\totan+\an:\totan:.5cm); 
          \fi
        \fi   
        \fill[\col,draw=black] (\totan:0.5cm)--(\totan:1cm)  arc(\totan:\totan+\an:1cm)
                                     --(\totan+\an:0.5cm) arc(\totan+\an:\totan :0.5cm);     
       \pgfmathsetmacro\finan{\totan+\an}
       \ifdim 180pt<\finan pt 
         \ifdim 180pt=\bottoman pt
            \shadedraw[left color=\col!50!\piechartthreedmixcolor,
                       right color=\col!25!\piechartthreedmixcolor,
                       draw=black,very thin] (180:1cm) -- ++(0,-3mm) arc (180:\totan+\an:1cm) 
                                                       -- ++(0,3mm)  arc (\totan+\an:180:1cm);
            \pgfmathsetmacro\bottoman{0}
            \global\let\bottoman\bottoman
            \else
            \shadedraw[left color=\col!50!\piechartthreedmixcolor,
                       right color=\col!25!\piechartthreedmixcolor,
                       draw=black,very thin](\totan:1cm)-- ++(0,-3mm) arc(\totan:\totan+\an:1cm)
                                                        -- ++(0,3mm)  arc(\totan+\an:\totan:1cm); 
          \fi
        \fi
        \pgfmathsetmacro\totan{\totan+\an}  \global\let\totan\totan 
       } 
    \end{scope}
    \draw[thin,black](0,0) circle (0.5cm);
   \end{scope}  
\end{scope}
}



%Tikz
\tikzset{
    %Hexagons
    hexagon/.style n args={3}{double arrow, double arrow head extend=0cm, inner sep=3pt, draw=#1, fill=#2, text=#3, thick},
    hexagon/.default={black}{gray!20}{black},
    hexagonOne/.style={hexagon={#1}{#1!30}{#1}},
    hexagonTwo/.style 2 args={hexagon={#1}{#1!#2}{#1}},
    hexagonThree/.style n args={3}{hexagon={#1}{#1!#2}{#3}},
    hexagonShade/.style n args={3}{double arrow, double arrow head extend=0cm, inner sep=3pt, thick, draw=#1, left color=#2, right color=#3},
    %General text element
    Shape/.style n args={3}{draw=#1, fill=#2, text=#3},
    genShape/.style 2 args={#1, inner sep=3pt, draw=#2, fill=#2!20, text=#2, thick},
    halo/.style={preaction={draw, #1, line width=7, -}},
}

%My commands
\newcommand<>{\tc}[2]{\textcolor#3{#1}{#2}}
\newcommand{\myCopyright}{\raisebox{-7pt}{\large\copyright}}
\makeatletter
    \newlength{\textsize}
    \setlength{\textsize}{\f@size pt}
    \newcommand{\emoji}[2][\;]{#1\raisebox{-0.15\textsize}{\includegraphics[height=\f@size pt]{#2}}}
\makeatother
\newcommand{\yes}[1][ForestGreen]{\textcolor{#1}{\ding{51}}}
\newcommand{\no}[1][OrangeRed]{\textcolor{#1}{\ding{55}}}

\makeatletter
\patchcmd{\beamer@sectionintoc}
  {\vfill}
  {\vskip\itemsep}
  {}
  {}
\makeatother

\mode<presentation>
{
    \usetheme{Z02}
    \defbeamertemplate{footline}{Empty}{}
    \setbeamersize{text margin left=5mm,text margin right=5mm}
}

%\pgfkeys{/beamer/overlayarea/show frame}

%Append code to put more information on titlepage
\appto\titlepage{%
    \begin{tikzpicture}[remember picture, overlay]
        \node[anchor=south, inner ysep=4pt] at ($(current page.south)!0.58!(current page.south east)+(0mm,1mm)$)
        {\includegraphics[width=13mm]{LogoPUNCH}};
        \node[anchor=south, inner ysep=0pt] (license) at ($(current page.south)+(0mm,4mm)$)
        {\includegraphics[width=10mm]{CC-BY}};
        \node[anchor=south, font=\small, text=BGLIGHT]
        (Axel) at (license.north) {\href{https://github.com/AxelKrypton}{\faGithub\,AxelKrypton}};
    \end{tikzpicture}
}

%===============================================================%
\title{Clean testing}
\subtitle{Good practices in general coding}
\author{Alessandro Sciarra}
\institute{Z02~--~Software Development Center\\[1mm] {PUNCH Young Academy -- PUNCH4NFDI}}
\date{24 May 2024}
\titlegraphic{\includegraphics[width=16mm]{LogoCRC}}
\titlepagelogo{\includegraphics[width=26mm]{LogoGoethe}}
%===============================================================%

\AtBeginSection[] % <- Empty optional argument, do nothing for \section*
{
    \begin{frame}[plain, noframenumbering]{}
         \sectionpage
    \end{frame}
}
\begin{document}
%===================================================================================================%
\begin{frame}[plain,noframenumbering]
    \titlepage
\end{frame}
%~~~~~~~~~~~~~~~~~~~~~~~~~~~~~~~~~~~~~~~~~~~~%
\begin{frame}[plain,noframenumbering]{Prelude}
    \vspace{-5mm}
    \begin{varblock}{block example}[0.8\textwidth]{Leave me a feedback}
        Was the talk clear, useful, inspiring?\\
        Any general comment?\\
        Drop me an email! $\;\to\;$
        \tc{PP}{\href{mailto:sciarra@th.physik.uni-frankfurt.de}{{\small\Letter}$\;$Alessandro}}
    \end{varblock}
    \vspace{1mm}
    \begin{varblock}{block alerted}[0.9\textwidth]{Disclaimer}
        Slides are quite full of text for later reading. \\
        Most example are oversimplified to fit on a slide.\\
        Tests shown before the TDD section are probably bad.\\[3mm]
        \textbf{\large Let's discuss!}
    \end{varblock}
\end{frame}
%~~~~~~~~~~~~~~~~~~~~~~~~~~~~~~~~~~~~~~~~~~~~%
\begin{frame}[plain,noframenumbering]{Outline}
    \vspace{-3mm}
    \hspace*{1cm}
    \begin{minipage}[0.65\textheight]{\textwidth}
        \tableofcontents
    \end{minipage}
\end{frame}
%===================================================================================================%


%===================================================================================================%
\section{A short recap about clean code}
%~~~~~~~~~~~~~~~~~~~~~~~~~~~~~~~~~~~~~~~~~~~~%
\begin{frame}{What we explored in the clean code talk}
    \begin{itemize}
        \item Take advantage of your IDE
        \item Use meaningful names
        \item Limit comments and care about formatting
        \item Function and classes should have a single responsibility (SRP)
        \item Integration Operation Separation Principle (IOSP)
        \item Don't repeat yourself (DRY)
        \item Keep it simple, stupid (KISS)
        \item Beware of optimisation
        \item The boyscout rule
        \item Keep improving
    \end{itemize}
    \begin{tikzpicture}[remember picture, overlay]
        \node[anchor=south east, inner sep=5mm] (fig) at (current page.south east)
            {\includegraphics[width=0.5\textwidth]{Bifurcation}};
        \node[text=BGLIGHT, font=\Large\bfseries] at ($(fig.center)!0.3!(fig.south)+(0.05,0)$)
            {
                \begin{tabular}{c}
                    Develop\\[-2mm]
                    your code sense\\
                \end{tabular}
            };
    \end{tikzpicture}
\end{frame}
%===================================================================================================%


%===================================================================================================%
\section[Why testing?]{Why (automated) testing?}
%~~~~~~~~~~~~~~~~~~~~~~~~~~~~~~~~~~~~~~~~~~~~%
\AlertFrame{A mind-blowing changing point!}
%~~~~~~~~~~~~~~~~~~~~~~~~~~~~~~~~~~~~~~~~~~~~%
\begin{frame}{Lots of software should rather work...}{In our society you cannot spend 60 seconds without interacting with a software system}
    \vspace{-2mm}
    \begin{enumerate}
        \item \parbox[t]{0.5\textwidth}{In modern car there is software that controls the steering wheel\ldots}
              \qquad\raisebox{-.5\height}{\includegraphics[width=22mm]{Lane_assistant_symbol}}
        \item Think of medical devices, aeroplanes, lifts, etc.
        \item Aren't you upset when an app on your phone stalls or crashes?!
    \end{enumerate}
    \vspace{2mm}
    \begin{center}
        \includegraphics[width=0.3\textwidth]{Airplane}
        \includegraphics[width=0.3\textwidth]{Lift}
        \includegraphics[width=0.3\textwidth]{Dentist_tools}
    \end{center}
\end{frame}
%~~~~~~~~~~~~~~~~~~~~~~~~~~~~~~~~~~~~~~~~~~~~%
\begin{frame}{An old, simple fact}{Taken from Casper Jones, «Program quality and programmer productivity» -- 1977}
    \vspace{-2mm}
    \begin{center}
        \includegraphics[width=0.82\textheight]{Bugs_per_line}
    \end{center}
\end{frame}
%~~~~~~~~~~~~~~~~~~~~~~~~~~~~~~~~~~~~~~~~~~~~%
\newcommand{\FermatExample}[4][3]{$#2^#1+#3^#1\neq#4^#1$}
\begin{frame}{The unfortunate tendency in many cases}
    \vspace{-3mm}
    \begin{itemize}
        \item \alert{It gives reasonable results $\;\Rightarrow\;$ it works!} \emoji{Scream}
        \item Manual (and possibly not systematic) testing is unfortunately still common practice in our community!
        \item Even if done in a systematic way, you will never run your tests enough!
    \end{itemize}
    \begin{overlayarea}{\textwidth}{0.4\textheight}
        \begin{onlyenv}<1-2>
            \vspace{3mm}
            \begin{varblock}{example}[\textwidth]{An remarkable story about manual testing}
                I have been using an own manually tested ``agenda-script'' to track my working time \alert{\textbf{for months}} and then one day it suddenly crashed.\\
                Where was the problem?
                \uncover<2>{\PB{I assumed \textbf{every} day has 24 hours\ldots\emoji{Innocent}}}
            \end{varblock}
        \end{onlyenv}
        \begin{onlyenv}<3>
            \begin{varblock}{quote}[0.7\textwidth]{Theorem:}
                Prove that, with $n,x,y,z\in \mathbb{N}>0$ and $n>2$, the equation $x^n+y^n =z^n$ has no solutions.
            \end{varblock}
            \begin{varblock}{quote}[0.8\textwidth]{Would you accept this as proof or even argument?}
                \small
                \FermatExample{2}{3}{4}\; and\;\FermatExample[4]{1}{5}{6}\; and\;\FermatExample[5]{4}{2}{5}
            \end{varblock}
        \end{onlyenv}
        \begin{onlyenv}<4>
            \vspace{3mm}
            \begin{varblock}{alert}[0.7\textwidth]{Now, software testing is not like mathematics...}
                \large ...but would you accept a physics fact \textbf{without rigorous evidence}?
            \end{varblock}
        \end{onlyenv}
    \end{overlayarea}
\end{frame}
%~~~~~~~~~~~~~~~~~~~~~~~~~~~~~~~~~~~~~~~~~~~~%
\begin{frame}{Consider one of your programs. Does it work?}
    \vspace{-5mm}
    \begin{varblock}{example}[0.8\textwidth]{How do you prove that a software works?}
        \uncover<2->{
            You don't.
            You try proving it doesn't work\ldots\ and fail!
        }
    \end{varblock}
    \begin{uncoverenv}<3>
        \begin{varblock}{}[\textwidth]{The scientific method}
            \begin{enumerate}
                \item Make an hypothesis $\PB{\;\mathcal{C}\;}\to\;$  «My code is correct»
                \item Tests attempt to show $\PB{\;!\mathcal{C}\;}$
                \item Confidence in $\;\PB{\mathcal{C}\;}$ tracks thoroughness of tests
            \end{enumerate}
        \end{varblock}
        \vspace{-2mm}
        \begin{varblock}{}[\textwidth]{}
            \large
            Testing code is science (cf. Popper's falsification principle)
        \end{varblock}
    \end{uncoverenv}
\end{frame}
%~~~~~~~~~~~~~~~~~~~~~~~~~~~~~~~~~~~~~~~~~~~~%
\begin{frame}{Uncle Bob's magic button}
    \begin{varblock}{quote}[1.02\textwidth]{Fearless competence}<2->
        Are you afraid to touch the code?
        Are you afraid to improve the code?
        Do you have that little echo in your  ear saying -- Ah, if it ain't broken, don't fix it?
        I expect you to improve the code, all of the time, fearless competence.
        How do you get fearless competence?
    \end{varblock}
    \vspace{-3mm}
    \begin{center}
        \begin{tikzpicture}
            \node (button) {\includegraphics[width=0.25\textwidth]{Magic_button}};
            \node[right = of button] (Y) {\includegraphics[width=12mm]{Yes}};
            \node[left = of button] (N) {\includegraphics[width=12mm]{X}};
            \path[to, thick] (button) edge[out=0, in=180] (Y);
            \path[to, thick] (button) edge[out=180, in=0] (N);
        \end{tikzpicture}
    \end{center}
    \vspace{-3mm}
    \begin{varblock}{quote}[0.95\textwidth]{}[Robert C. Martin]<2->
        Tests! What if I had a button, a little button I could push on my keyboard and some lights would blink and science-fiction sounds would come out of my laptop for a few minutes and then a little green light would light up and that green light told me that the system worked.
        \alert{And I believed it.}
    \end{varblock}
\end{frame}
%~~~~~~~~~~~~~~~~~~~~~~~~~~~~~~~~~~~~~~~~~~~~%
\begin{frame}{The correct mental approach to testing}
    \begin{itemize}
        \setlength{\itemindent}{-2mm}
        \item Testing is part of the natural evolution of the code (not an overhead)
        \item Test code has to fulfil the same clean code rules as production code
        \item \alert{Tests are the key to be free to change (improve) the production code!}
        \item Keep balance between adding new code, writing tests and refactoring
    \end{itemize}
    \vspace{3mm}
    \begin{center}
        \resizebox{!}{0.32\textheight}{
            \begin{tikzpicture}
                \tikzstyle{Pie}=[font=\Large]
                \tikzstyle{PieNum}=[font=\Large\bfseries\ttfamily]
                \piechartthreed[scale=0.8,
                background color=white,
                mix color= white]{120/PS!50!LimeGreen,120/PP!50!white,120/PT}
                \foreach \i in {1,...,3} { \fill (pc \i) ellipse (0.1 and 0.05);}
                \draw[darkgray] (pc 1)  -- ++(7,0) coordinate (s1)
                node[Pie, anchor=south east] {Testing}
                node[PieNum, anchor=north east] {\tc{PS}{33\%}};
                \draw[darkgray] (pc 2)  -- ++(-1,1) -- ++(-5,0) coordinate (s2)
                node[Pie, anchor=south west] {Developing}
                node[PieNum, anchor=north west] {\tc{PP}{33\%}};
                \draw[darkgray] (pc 3)  -- (pc 3 -| s2)
                node[Pie, anchor=south west] {Refactoring}
                node[PieNum, anchor=north west] {\tc{PT}{33\%}};
            \end{tikzpicture}
        }
    \end{center}
    \FrameRemark{Slide taken from the \textsl{Clean Code} talk.}
    \begin{tikzpicture}[overlay, remember picture]
        \node[starburst, starburst points=27, starburst point height=16mm,
              minimum width=12cm, minimum height=5cm, Shape={red}{yellow!50}{red},
              line width=1.5pt, font=\Huge, rotate=10, visible on=<2>]
          at (current page.center) {Today much more!};
    \end{tikzpicture}
\end{frame}
%===================================================================================================%


%===================================================================================================%
\section{Different types of testing}
%~~~~~~~~~~~~~~~~~~~~~~~~~~~~~~~~~~~~~~~~~~~~%
\begin{frame}{Software testing}{Test code is just as important as production code}
    \vspace{-7mm}
    \begin{columns}[t]
        \begin{column}{0.45\textwidth}
            \begin{itemize}
                \item \PP{Testing levels}
                \begin{itemize}
                    \item \PS{Unit testing}
                    \item \PS{Integration testing}
                    \item System testing
                    \item Acceptance testing
                \end{itemize}
                \vspace{6mm}
                \centering
                \includegraphics[height=28mm]{SoftwareTesting.jpeg}\hfill
                \includegraphics[height=28mm]{TDD.jpeg}
            \end{itemize}
        \end{column}
        \begin{column}{0.55\textwidth}
            \begin{itemize}
                \item \PP{Testing types and techniques}
                \begin{itemize}
                    \item \textcolor{bg!80!normal text.fg}{Manual testing}
                    \item \PS{Automated testing}
                    \item Continuous testing
                    \item Functional testing
                    \item Mutation testing
                    \item Fuzzy testing
                    \item \ldots
                \end{itemize}
                \vspace{2mm}
                \item \small \PB{Plenty of videos/resources on the web}
            \end{itemize}
            \begin{itemize}
                \scriptsize\setlength{\itemsep}{0mm}
                \item[]
                \URL{https://www.youtube.com/watch?v=u5senBJUkPc}{~CppCon2015 -- All your tests are terrible\ldots}
                \item[]
                \URL{https://www.youtube.com/watch?v=FjwayiHNI1w}{~CppCon2020 -- The science of unit tests}
                \item[]
                \URL{https://www.youtube.com/watch?v=_OHE33s7EKw}{~CppCon2020 -- Back to Basics: Unit Tests}
                \item[]
                \URL{https://www.youtube.com/watch?v=EZ05e7EMOLM}{~TDD, Where Did It All Go Wrong}
                \item[]
                \URL{https://www.youtube.com/watch?v=a6oP24CSdUg}{~TDD for those who don't need it}
                \item[] \ldots
            \end{itemize}
        \end{column}
    \end{columns}
    \FrameRemark{The videos linked here above (as well as the references therein) were mainly used to prepare this talk.}
\end{frame}
%~~~~~~~~~~~~~~~~~~~~~~~~~~~~~~~~~~~~~~~~~~~~%
%Tikz group of commands to make a pyramid layer
\newcommand{\pyramidLayer}[6][]{%
    \pgfmathsetmacro{\Thickness}{#3}
    \pgfmathsetmacro{\Long}{#4}
    \pgfmathsetmacro{\Short}{#4-0.5*#3}
    \path coordinate (middleUpperRight) at ($(#2)+(0,\Thickness)+(30:0.5*\Short)$)
    coordinate (middleLowerRight) at ($(#2)+(30:0.5*\Long)$)
    coordinate (upperRight) at ($(#2)+(0,\Thickness)+(30:\Short)$)
    coordinate (lowerRight) at ($(#2)+(30:\Long)$);
    \draw[fill=#5!50] ($(#2)+(0,\Thickness)+(30:\Short)$)  -- ++(210:\Short) -- ++(270:\Thickness) -- ++(30:\Long) -- cycle;
    \draw[fill=#5!20] ($(#2)+(0,\Thickness)+(150:\Short)$)  -- ++(330:\Short) -- ++(270:\Thickness) -- ++(150:\Long) -- cycle;
    \draw[fill=#5]    ($(#2)+(0,\Thickness)$) -- ++(30:\Short) -- ++(150:\Short) -- ++(210:\Short) -- cycle;
    \node[font=\small] at ($(middleUpperRight)!0.5!(middleLowerRight)$) {\rotatebox{30}{#6}};
    \node[text=#5, anchor=west] at ($(upperRight)!0.5!(lowerRight)+(0.5,0)$) {#1};
}
\begin{frame}{The tests pyramid}
    \begin{center}
        \resizebox{!}{0.8\textheight}{
            \begin{tikzpicture}
                \path coordinate (A) at (0,0)
                      coordinate (B) at (0,2)
                      coordinate (C) at (0,4)
                      coordinate (D) at (0,6)
                      coordinate (O) at ($(A)+(150:5)-(0.5,0)$);
                \draw (O) -- ++(30:2); %Just draw axis behind pyramid
                %Draw Pyramid
                \pyramidLayer{A}{1.5}{5}{red}{Unit};
                \pyramidLayer{B}{1.5}{4}{orange!80!yellow}{Integration};
                \pyramidLayer{C}{1.5}{3}{PS}{System};
                \pyramidLayer{D}{1.5}{2}{PP}{Acceptance};
                \draw[to] (O) -- ++(0,6) node[right, anchor=north west] {
                    \begin{tabular}{@{}l}
                        Development cost\\
                        Execution time\\
                    \end{tabular}
                };
                \draw[to] (O) -- ++(330:6) node[right] {Number of tests};
            \end{tikzpicture}
        }
    \end{center}
\end{frame}
%~~~~~~~~~~~~~~~~~~~~~~~~~~~~~~~~~~~~~~~~~~~~%
\begin{frame}{Unit tests are not enough!}
    \begin{varblock}{}[0.9\textwidth]{Integration tests}
        \PB{Once you are sure that all components work correctly alone, be sure they work as expected together!}
    \end{varblock}
    \vspace{3mm}
    \begin{overlayarea}{\textwidth}{0.8\textheight}
        \begin{center}
            \begin{onlyenv}<1>
                \includegraphics[height=4cm]{UnitTestVSIntegrationTestA}$\quad$
                \includegraphics[height=4cm]{UnitTestVSIntegrationTestB}\\
            \end{onlyenv}
            \begin{onlyenv}<2>
                \includegraphics[height=35mm]{MarsClimateOrbiter}\\
                {
                    \scriptsize
                    \URL[PB]{http://www.vitalstatistics.info/uploads/mars\%20climate\%20orbiter.jpg}{Mars Climate Orbiter}
                    lost in 1999 -- By NASA/JPL/Corby Waste -- \myCopyright\ Public Domain
                }
            \end{onlyenv}
        \end{center}
    \end{overlayarea}
    \FrameRemark{I already mentioned this in the \textsl{Clean code} talk.}
\end{frame}
%~~~~~~~~~~~~~~~~~~~~~~~~~~~~~~~~~~~~~~~~~~~~%
\begin{frame}{And once integration tests pass?}
    \vspace{-2mm}
    \begin{varblock*}{quote}[0.86\textwidth]{System testing}[Wikipedia]
        System testing tests a completely integrated system to verify that the system meets its requirements.
    \end{varblock*}
    \vspace{-3mm}
    \begin{varblock*}{quote}[0.86\textwidth]{Acceptance testing}[Standard Glossary of Terms used in Software Testing]
        Formal testing with respect to user needs, requirements, and business processes conducted to determine whether a system satisfies the acceptance criteria and to enable the user, customers or other authorized entity to determine whether to accept the system.\\[-2mm]
    \end{varblock*}
    \vspace{-3mm}
    \begin{varblock}{example}[\textwidth]{In smaller projects}
        You might not need all categories, but, if you drop any, do it consciously!
    \end{varblock}
\end{frame}
%~~~~~~~~~~~~~~~~~~~~~~~~~~~~~~~~~~~~~~~~~~~~%
\begin{frame}{A technique for system/acceptance tests}
    \begin{varblock*}{quote}[0.86\textwidth]{Functional tests}[Wikipedia]
        They tend to answer the question of ``can the user do this'' or ``does this particular feature work''.
    \end{varblock*}
    \begin{center}
        \begin{tikzpicture}[thick, node distance=15mm]
            \node[cylinder, cylinder uses custom fill, cylinder end fill=PP!50,
                  cylinder body fill=PP!25, aspect=0.5, draw=PB, minimum height=15mm,
                  minimum width=12mm] (SW) {Software};
            \begin{scope}[every node/.style={rectangle, rounded corners}]
                \node[left = of SW,  draw=PS, fill=PS!20, yshift=+8mm] (IN) {Given input};
                \node[right = of SW, draw=PT, fill=PT!20, yshift=-8mm] (OUT) {Given output};
            \end{scope}
            \path[to, thin] (IN) edge[out=0, in=180] (SW);
            \path[to, thin] ($(SW.before top)!0.5!(SW.after top)$) edge[out=0, in=180] (OUT);
        \end{tikzpicture}
    \end{center}
    \begin{itemize}
        \item Consider to implement these tests in an external script!
        \item \alert{This is a good starting point to work with legacy code without tests!}
    \end{itemize}
\end{frame}
%===================================================================================================%


%===================================================================================================%
\section{Tests goals -- properties -- good principles}
%~~~~~~~~~~~~~~~~~~~~~~~~~~~~~~~~~~~~~~~~~~~~%
\begin{frame}{General idea and overview}
    \begin{overlayarea}{\textwidth}{0.7\textheight}
        \begin{center}
            \large
            \only<1-2>{There is slightly different advice around, pick your favourite:}
            \only<3->{There is one not-negotiable \textbf{requirement of tests}, though:}
        \end{center}
        \begin{varblock}{example}[0.8\textwidth]{Write F.I.R.S.T tests first}<only@1>
            \begin{description}[Self-validating:]
                \item[\PQ{F}ast:] To be run frequently
                \item[\PQ{I}ndependent:] To make failures meaningful
                \item[\PQ{R}epeatable:] To be run everywhere with same outcome
                \item[\PQ{S}elf-validating:] To be automatised
                \item[\PQ{T}imely:] To make production code testable
            \end{description}
        \end{varblock}
        \begin{varblock}{example}[\textwidth]{Google's perspective}<only@2>
            \begin{description}[Self-validating:]
                \item[Readable:] Tests are correct by inspection
                \item[Correct:] Do not rely on known bugs and test real scenarios
                \item[Complete:] Test most edge cases
                \item[Documenting:] Demonstration of how the API works
                \item[Resilient:] Repeatable, independent, hard to break, hermetic, etc.
            \end{description}
        \end{varblock}
        \begin{tikzpicture}[overlay, remember picture]
            \node[starburst, starburst points=21, starburst point height=15mm,
                  minimum width=10cm, minimum height=6cm, Shape={red}{yellow!50}{red},
                  line width=1.5pt, font=\Huge\bfseries, visible on=<4>]
                at ($(current page.center)-(0,1)$) {Existence};
        \end{tikzpicture}
    \end{overlayarea}
\end{frame}
%~~~~~~~~~~~~~~~~~~~~~~~~~~~~~~~~~~~~~~~~~~~~%
\begin{frame}[fragile]{A small interlude: Unit test frameworks for C++}{Essential for self-validation and automation!}
    \vspace{2mm}
    \begin{overlayarea}{\textwidth}{0.8\textheight}
        \begin{enumerate}[<only@+->]
            \setlength{\itemindent}{3mm}
            \item \tikzmark{Boost} Boost Test Library\hfill
                  \parbox{0.4\textwidth}{\footnotesize\URL[PB]{https://www.boost.org/doc/libs/1_82_0/libs/test/doc/html/index.html}{Documentation}}
            \item \tikzmark{Google} GoogleTest \uncover<+->{\;\textcolor{bg!60!normal text.fg}{or}\; Catch2}\hfill
                  \parbox{0.4\textwidth}{
                      \footnotesize
                      \URL[PB]{https://google.github.io/googletest/}{User's guide} \uncover<.->{or
                      \URL[PB]{https://github.com/catchorg/Catch2}{GitHub project}}
                  }
            \item \tikzmark{Vir} Vir's Unit Test Framework\hfill
                  \parbox{0.4\textwidth}{\footnotesize\URL[PB]{https://github.com/mattkretz/virtest}{GitHub repository}}
            \item[] {\small ...and so many mores!}\hfill
                  \parbox{0.4\textwidth}{
                      \footnotesize
                      \URL[PB]{https://en.wikipedia.org/wiki/List_of_unit_testing_frameworks\#C++}{List on Wikipedia}
                  }
        \end{enumerate}
        \begin{varblock}{example}[0.9\textwidth]{}<only@1>
            \begin{lstlisting}[style=MyCpp]
                #define BOOST_TEST_MODULE My Test
                #include |+<boost/test/included/unit_test.hpp>+|

                BOOST_AUTO_TEST_CASE(first_test)
                {
                    int i = 1;
                    BOOST_TEST(i);
                    BOOST_TEST(i == 2);
                }
            \end{lstlisting}
            \begin{lstlisting}[style=MyBash]
                |+Running 1 test case...
                test_file.cpp(8): error: in "first_test": check i == 2 has failed [1 != 2]

                *** 1 failure is detected in the test module "My Test"+|
            \end{lstlisting}
        \end{varblock}
        \begin{varblock}{example}[0.9\textwidth]{}<only@2>
            \begin{lstlisting}[style=MyCpp]
                #include |+<gtest/gtest.h>+|
                #include |+"src/factorial.h"+|

                TEST(FactorialTest, HandlesZeroInput) {
                    EXPECT_EQ(factorial(0), 1);
                }

                TEST(FactorialTest, HandlesPositiveInput) {
                    EXPECT_EQ(factorial(1), 1);
                }

                int main(int argc, char **argv) {
                    ::testing::InitGoogleTest(&argc, argv);
                    return RUN_ALL_TESTS();
                }
            \end{lstlisting}
        \end{varblock}
        \begin{varblock}{example}[0.9\textwidth]{}<only@3>
            \begin{lstlisting}[style=MyCpp]
                #include |+<catch2/catch_test_macros.hpp>+|

                unsigned int factorial(unsigned int number)
                {
                    return number > 1 ? factorial(number - 1) * number : 1;
                }

                TEST_CASE("Factorials are computed", "[factorial]") {
                    REQUIRE(factorial(0)  == 1);
                    REQUIRE(factorial(1)  == 1);
                    REQUIRE(factorial(3)  == 6);
                    REQUIRE(factorial(10) == 3628800);
                }
            \end{lstlisting}
        \end{varblock}
        \begin{varblock}{example}[0.9\textwidth]{}<only@4>
            \begin{lstlisting}[style=MyCpp]
                #include |+<vir/test.h>+|

                TEST(test_name) {
                    int test = 3;
                    COMPARE(test, 2) << "more " << "details";
                    //COMPARE(test, 2).on_failure("more ", "details");
                }
            \end{lstlisting}
            \begin{lstlisting}[style=MyBash]
                @|\tikzmark{corner}|@|+FAIL:  at tests/testfile.cpp:5 (0x40451f):
                FAIL:  test (3) == 2 (2) -> false more details
                FAIL:  test_name

                Testing done. 0 tests passed. 1 tests failed.+|
            \end{lstlisting}
        \end{varblock}
        \vspace{1mm}
        \begin{varblock*}{alert}[0.8\textwidth]{Not all have the same weight and functionality}<only@5>
            Pick your favourite depending on your needs and project:
            \vspace{2mm}
            \begin{enumerate}
                \small
                \item \tikzmark{BoostInfo} Natural choice if you already depend on Boost
                \item \tikzmark{GoogleInfo} More complex but also more powerful
                \item \tikzmark{VirInfo} Very simple and lightweight
            \end{enumerate}
        \end{varblock*}
        \begin{tikzpicture}[remember picture, overlay]
            % virtest output symbol
            \draw[thick, visible on=<4>] ($(corner)-(0.335\textwidth,0.13)$) -- ++(-0.12,0) -- ++(0,-0.68) -- ++(0.12,0);
            % connecting lines of enumerate environments
            \begin{scope}[visible on=<5>]
                \draw[thin, to] ($(Boost)-(0.68,0)$)  -- ++(-0.45,0) |- ($(BoostInfo)-(0.45\textwidth,0.13)$);
                \draw[thin, to] ($(Google)-(0.68,0)$) -- ++(-0.30,0) |- ($(GoogleInfo)-(0.45\textwidth,0.13)$);
                \draw[thin, to] ($(Vir)-(0.68,0)$)    -- ++(-0.15,0) |- ($(VirInfo)-(0.45\textwidth,0.13)$);
            \end{scope}
        \end{tikzpicture}
    \end{overlayarea}
    \FrameRemark<5>{If you use a different language, you can find at the Wikipedia link lists for most programming languages.}
\end{frame}
%~~~~~~~~~~~~~~~~~~~~~~~~~~~~~~~~~~~~~~~~~~~~%
\begin{frame}[fragile]{Having tests is not enough, you must trust them!}{Few words about tests correctness}
    \begin{overlayarea}{\textwidth}{0.7\textheight}
        \begin{itemize}
            \item Never rely on known bugs
            \item<only@4-> Never mock the code that you are testing
        \end{itemize}
        \begin{varblock}{alert}[0.6\textwidth]{Having tests passing is simple...}<only@1-2>
            \begin{uncoverenv}<2->
                \begin{lstlisting}[style=MyCpp]
                    int cube(int x) {
                        //TODO: To be implemented
                        return 0;
                    }

                    TEST(cube_test) {
                        VERIFY(0 == cube(2));
                        VERIFY(0 == cube(3));
                        VERIFY(0 == cube(42));
                    }
                \end{lstlisting}
            \end{uncoverenv}
        \end{varblock}
        \begin{varblock}{example}[0.6\textwidth]{...better to have them failing!}<only@3>
            \begin{lstlisting}[style=MyCpp]
                int cube(int x) {
                    //TODO: To be implemented
                    return 0;
                }

                TEST(cube_test) {
                    VERIFY(8 == cube(2));
                    VERIFY(27 == cube(3));
                    VERIFY(74088 == cube(42));
                }
            \end{lstlisting}
        \end{varblock}
        \begin{varblock}{alert}[0.8\textwidth]{This is probably useless}<only@4>
            \begin{lstlisting}[style=MyCpp]
                class MockCoffeeMachine : public CoffeeMachine {
                    //For simplicity we assume always in temperature
                    double warmUp() override {
                        return constants::brewingT;
                    }
                }

                BOOST_AUTO_TEST_CASE(warmUp_test) {
                    MockCoffeeMachine machine{};
                    BOOST_REQUIRE_CLOSE(machine.warmUp(),
                                        constants::brewingT, 0.1);
                }
            \end{lstlisting}
        \end{varblock}
        \vspace{3mm}
        \begin{varblock}{}[0.85\textwidth]{Think of a clear reaction procedure in your team}<only@5->
            What should happen in your project when a test fails?\\
            ...and when a bug that was not caught by tests is found?\\
            $\downarrow$\\
            \PB{Never miss any chance to improve your tests!}
        \end{varblock}
    \end{overlayarea}
\end{frame}
%~~~~~~~~~~~~~~~~~~~~~~~~~~~~~~~~~~~~~~~~~~~~%
\begin{frame}{Tests readability}
    \begin{center}
        \begin{tikzpicture}[every node/.style={inner sep=0mm, text=BGLIGHT}]
            \node (fig) {\includegraphics[width=0.75\textwidth]{Balance}};
            \begin{DrawOnNode}{fig}
                \node[anchor=east] at (0.41,0.78) {
                    \begin{tabular}{c}
                        Avoid boilerplate\\
                        \footnotesize E.g. use functions for setup\\
                    \end{tabular}
                };
                \node[anchor=west] at (0.51,0.78) {
                    \begin{tabular}{c}
                        Give enough context\\
                        \footnotesize E.g. Not a single function for setup\\
                    \end{tabular}
                };
            \end{DrawOnNode}
        \end{tikzpicture}
    \end{center}
    \begin{varblock}{}[0.85\textwidth]{}
        \Large Remember: \PB{Tests are correct by inspection!}
    \end{varblock}
\end{frame}
%~~~~~~~~~~~~~~~~~~~~~~~~~~~~~~~~~~~~~~~~~~~~%
\begin{frame}[fragile]{Tests completeness: What should I test and how?}
    \begin{onlyenv}<1>
        \begin{itemize}
            \item Values  \\ \centerline{\includegraphics[scale=0.1]{TestingValue}}
            \item Actions \\ \centerline{\includegraphics[scale=0.12, clip, trim=0 4cm 0 4cm]{TestingAction}}
        \end{itemize}
        \vspace{3mm}
        \hfill\ldots{}and anything that makes sense!
    \end{onlyenv}
    \begin{varblock}{example}[0.98\textwidth]{A very nice example from Titus Winter}<only@2->
        \begin{overlayarea}{\textwidth}{0.55\textheight}
            \begin{uncoverenv}<3->
                \begin{lstlisting}[style=MyCpp]
                    int Factorial(int n) {
                        if (n == 1) return 1;
                        if (n == 5) return 120;
                        return -1; // TODO(goofus): figure this out.
                    }
                \end{lstlisting}
            \end{uncoverenv}
            \begin{onlyenv}<2-3>
                \begin{lstlisting}[style=MyCpp]
                    TEST(FactorialTest, BasicTests) {
                        EXPECT_EQ(1, Factorial(1));
                        EXPECT_EQ(120, Factorial(5));
                    }
                \end{lstlisting}
            \end{onlyenv}
            \begin{onlyenv}<4>
                \begin{lstlisting}[style=MyCpp]
                    TEST(FactorialTest, BasicTests) {
                        EXPECT_EQ(1, Factorial(1));
                        EXPECT_EQ(120, Factorial(5));
                        EXPECT_EQ(1, Factorial(0));
                        EXPECT_EQ(479001600, Factorial(12));
                        EXPECT_EQ(std::numeric_limits<int>::max() , Factorial(13));
                        EXPECT_EQ(1, Factorial(0));
                        EXPECT_EQ(std::numeric_limits<int>::max() , Factorial(-10));
                    }
                \end{lstlisting}
            \end{onlyenv}
        \end{overlayarea}
    \end{varblock}
    \begin{uncoverenv}<4>
        \begin{center}
            This will pretty much naturally emerge when doing TDD
        \end{center}
    \end{uncoverenv}
    \FrameRemark<1>{I already mentioned this in the \textsl{Clean code} talk.}
\end{frame}
%~~~~~~~~~~~~~~~~~~~~~~~~~~~~~~~~~~~~~~~~~~~~%
\begin{frame}{The hidden value of tests}{Few words about tests as API usage examples}
    \begin{overlayarea}{\textwidth}{0.68\textheight}
        \begin{center}
            \small
            Of course, a meme, do not take it (too) seriously, but...\\
            \vspace{3mm}
            \begin{uncoverenv}<2>
                \includegraphics[width=0.7\textwidth]{Documentation_meme}\\
                \vspace{3mm}
                ...good tests should show how to use your code!
            \end{uncoverenv}
        \end{center}
        \FrameRemark<2>{In this sense, it should be avoided to use implementation details in tests (cf. white box testing downsides).}
    \end{overlayarea}
\end{frame}
%~~~~~~~~~~~~~~~~~~~~~~~~~~~~~~~~~~~~~~~~~~~~%
\begin{frame}[fragile]{Tests resilience: Google advice}
    \begin{onlyenv}<1>
        \begin{itemize}
            \setlength{\itemsep}{2mm}
            \item No flaky tests $\to$ Tests should always give the same outcome
            \item No brittle tests $\to$ One failure should not trigger many failures
            \item Tests reference values should not come from the system under test!
            \item Changing tests order should never change the outcome
            \item Tests should be as hermetic as possible $\to$ No I/O, no network, etc.
            \item No deep dependence $\to$ Avoid any implicit assumption
        \end{itemize}
        \vspace{3mm}
        \begin{varblock}{example}[0.5\textwidth]{}
            \large\bfseries\PS{Let's see some examples!}
        \end{varblock}
    \end{onlyenv}
    \begin{onlyenv}<2-3>
        \begin{varblock}{alert}[\textwidth]{A sneaky flaky test: What happens if the test-system is overloaded?}
            \begin{lstlisting}[style=MyCpp]
                TEST(UpdaterTest, RunsFast) {
                    Updater updater;
                    updater.UpdateAsync();
                    SleepFor(Seconds(.5)); // Half a second should be *plenty*.
                    EXPECT_TRUE(updater.Updated());
                }
            \end{lstlisting}
        \end{varblock}
        \begin{varblock}{alert}[\textwidth]{An infamous brittle pattern}<uncover@3>
            \begin{lstlisting}[style=MyCpp]
                TEST(Logger, LogWasCalled) {
                    StartLogCapture();
                    EXPECT_TRUE(Frobber::Start());
                    EXPECT_THAT(
                      Logs(), Contains("file.cc:421: Opened file frobber.config" )
                    );@|\tikzmark{logger}|@
                }
            \end{lstlisting}
        \end{varblock}
        \begin{tikzpicture}[remember picture, overlay]
            \draw[red, thick, double, visible on=<3>]
                ($(logger)-(2.6,0)$) -- ++(1.15,0)
                +(0.18,0) -- ++(0.65,0)
                +(2.4,0) -- ++(4.8,0);
        \end{tikzpicture}
    \end{onlyenv}
    \begin{onlyenv}<4-5>
        \begin{varblock}{alert}[\textwidth]{Where are the reference values coming from?}
            \begin{lstlisting}[style=MyCpp]
                BOOST_AUTO_TEST_CASE(Dirac_M)
                {
                    const hmc_float refs[4] =
                        { 2610.3804893063798, 4356.332327032359,
                          2614.2685771909237, 4364.1408252701831 };
                    test_fermionmatrix<physics::fermionmatrix::Dslash>(refs);
                }
            \end{lstlisting}
        \end{varblock}
        \begin{varblock}{alert}[\textwidth]{A non-hermetic test: What if run twice in parallel?}<uncover@5>
            \begin{lstlisting}[style=MyCpp]
                TEST(Server, StorageTest) {
                    StorageServer* server = GetStorageServerHandle();
                    auto my_val = rand();
                    server->Store("testkey", my_val);
                    EXPECT_EQ(my_val, server->Load("testkey"));
                }
            \end{lstlisting}
        \end{varblock}
    \end{onlyenv}
    \begin{onlyenv}<6>
        \vspace{-3mm}
        \begin{varblock}{alert}[\textwidth]{Deep dependence example}<only@6>
            \begin{lstlisting}[style=MyCpp]
                class File {
                  public:
                    // ...
                    virtual bool StatWithOptions(Stat_t* stat, StatOptions opts) {
                        return this->Stat(stat); // In base class ignore options
                    }
                };
                TEST(Filesystem, FSUsage) {
                    // ... and call to StatWithOptions
                    EXPECT_CALL(file, Stat(_)).Times(1);
                }  // This relies on a call to base StatWithOptions!!
            \end{lstlisting}
        \end{varblock}
        \begin{varblock}{quote}[\textwidth]{The law of implicit interfaces (AKA Hyrum's law)}[Hyrum Wright]<only@6>
            Given enough users of your public interface, soon or later someone will start implicitly relying on your implementation details (speed, memory consumption, etc.).\\[-1.2ex]~
        \end{varblock}
    \end{onlyenv}
    \FrameRemark{This list and all GoogleTest examples are taken from Titus Winters' and Hyrum Wright's talk at CppCon2015.}
\end{frame}
%~~~~~~~~~~~~~~~~~~~~~~~~~~~~~~~~~~~~~~~~~~~~%
\begin{frame}{A slide about tests precision}
    \vspace{-3mm}
    \begin{onlyenv}<1>
        \begin{center}
            \includegraphics[width=0.7\textwidth]{Precision_meme}
        \end{center}
    \end{onlyenv}
    \begin{onlyenv}<2>
        \begin{itemize}
            \item We all know that machine precision is finite!
            \item Comparing floating point numbers has to be done carefully
            \item If you know a reference value exactly, put in the code more than the needed digits
                  \PB{$\;\to\pi=3.141592653589793238462643$}
            \item Be aware of possible compiler techniques, e.g.\ FMA \Remark{fused multiply-add}\\
                  \then This could make your test fail if you require too high precision
            \item Requiring too high precision makes tests brittle
            \item Requiring too low precision makes tests useless
            \item \alert{After all it is a judgement call!}
        \end{itemize}
    \end{onlyenv}
\end{frame}
%~~~~~~~~~~~~~~~~~~~~~~~~~~~~~~~~~~~~~~~~~~~~%
\begin{frame}{Few more words about reproducibility}
    \vspace{-5mm}
    \begin{varblock}{example}[0.96\textwidth]{How can I make my tests reproducible?}
        \begin{enumerate}
            \item Isolate them from non-deterministic aspects \Remark{e.g.\ network, database, random numbers}\\
                  \then This is not always possible!
            \item Measure unwanted environment effect and subtract it in tests
            \item Detect if test can be run in the present environment\\
                  \then If not, either fail or handle situation
            \item \alert{Handle intrinsic noise by statistics}\\
                  \Then{E.g.\ if random numbers are needed, measure the failure frequency and re-run the test a number of times large enough to consider a failure as such!}
        \end{enumerate}
    \end{varblock}
\end{frame}
%~~~~~~~~~~~~~~~~~~~~~~~~~~~~~~~~~~~~~~~~~~~~%
\begin{frame}{Tests accuracy and reacting to failures}
    \begin{center}
        \begin{tikzpicture}
            \coordinate (O) at (0,0);
            \begin{scope}[every node/.style={inner ysep=6mm, inner xsep=3mm, font=\huge, minimum width=46mm},
                          every label/.style={label distance=-3mm, font=\small, anchor=south, inner sep=0mm}]
                \node[anchor=south east, fill=PP!30, visible on=<2->] (M11) at (O) {
                    \begin{tabular}{c}
                        \emoji[]{Top}\\[-5mm]
                        {\normalsize Ship it!}\\[-9pt]
                    \end{tabular}
                };
                \node[anchor=south west, fill=ALERT, visible on=<3->] (M12) at (O) {
                    \begin{tabular}{c}
                        \emoji[]{Warning}\\[-5mm]
                        {\normalsize Add tests, fix code!}\\[-9pt]
                    \end{tabular}
                };
                \node[anchor=north east, fill=PT!10, visible on=<4->] (M21) at (O) {
                    \begin{tabular}{c}
                        \emoji[]{Sweat_smile}\\[-5mm]
                        {\normalsize False alarm, fix tests!}\\[-9pt]
                    \end{tabular}
                };
                \node[anchor=north west, fill=PS!30, visible on=<5->] (M22) at (O) {
                    \begin{tabular}{c}
                        \emoji[]{Tada}\\[-5mm]
                        {\normalsize Success story!}\\[-9pt]
                    \end{tabular}
                };
            \end{scope}
            \begin{scope}
                \draw (M11.south west) -- (M11.north west) -- (M12.north east) -- (M22.south east) --
                      (M21.south west) -- (M21.north west) -- (M12.south east)
                      (M11.north east) -- (M21.south east);
            \end{scope}
            \begin{scope}[node distance=2mm, every node/.style={inner sep=1mm, font=\Large}]
                \node[above = of M11] (CCyes) {\yes};
                \node[above = of M12] (CCno)  {\no};
                \node[left  = of M11, yshift=+3mm] (TPyes) {\yes};
                \node[left  = of M21, yshift=-3mm] (TPno)  {\no};
            \end{scope}
            \begin{scope}[every node/.style={font=\small}]
                \node at ($(CCyes.east)!0.5!(CCno.west)+(0,0)$) {Is code correct?};
                \node[rotate=90] at ($(TPyes.south)!0.5!(TPno.north)-(0,0)$) {Did tests pass?};
            \end{scope}
        \end{tikzpicture}
    \end{center}
    \FrameRemark{Tests accuracy means how much tests outcome matches code correctness.}
\end{frame}
%===================================================================================================%


%===================================================================================================%
\section{The curse of the \textbf{unit} word}
%~~~~~~~~~~~~~~~~~~~~~~~~~~~~~~~~~~~~~~~~~~~~%
\begin{frame}{What is a \textbf{unit} in unit testing?}
    \begin{overlayarea}{\textwidth}{0.38\textheight}
        \begin{onlyenv}<1>
            \begin{varblock}{quote}[0.7\textwidth]{}[Wikipedia]
                A group of related functions often called a module.
            \end{varblock}
            \begin{varblock}{quote}[0.7\textwidth]{}[Ian Cooper]
                A unit test should test a requirement, only test against the stable contract of the API.
            \end{varblock}
        \end{onlyenv}
        \begin{onlyenv}<2>
            \vspace{-4mm}
            \begin{center}
                \includegraphics[width=0.7\textwidth]{Fragile_test_problem}
            \end{center}
        \end{onlyenv}
    \end{overlayarea}
    \begin{varblock}{alert}[0.98\textwidth]{A unit is not necessarily a class or a function}
        This is possibly one of the most misunderstood concepts in software engineering in the last decade.
        Sometimes a unit is a class or a function, but there should be a good reason for such a correspondence.
        \textbf{More often a unit should be (the public API of) a module.}
    \end{varblock}
\end{frame}
%~~~~~~~~~~~~~~~~~~~~~~~~~~~~~~~~~~~~~~~~~~~~%
\AlertFrame{
    \begin{tabular}{c}
        Wait, what's wrong\\ in testing each single class?
    \end{tabular}
}
%~~~~~~~~~~~~~~~~~~~~~~~~~~~~~~~~~~~~~~~~~~~~%
\begin{frame}{It depends}
    \vspace{-3mm}
    \begin{varblock}{quote}[0.8\textwidth]{TDD, Where Did It All Go Wrong}[Ian Cooper]
        A couple of times we were very proud of the fact that we produced some code that went to production with almost zero defects.
        [...]
        I looked at this test suite.
        One things I noticed was that \alert{we had about three times as much test code as production code}.
    \end{varblock}
    \begin{itemize}
        \item Some classes are simply implementation details\\
              \then You would like to freely change them
        \item Imagine a module with few classes working together \Remark{in your implementation, today}\\
              \then Does it make sense to mock all but one several times?
        \item If a class happens to be the API of a module, probably that's a unit
    \end{itemize}
\end{frame}
%===================================================================================================%


%===================================================================================================%
\section[White- and black-box testing!?]{The weaknesses of white- and black-box testing}
%~~~~~~~~~~~~~~~~~~~~~~~~~~~~~~~~~~~~~~~~~~~~%
\AlertFrame{
    \begin{tabular}{c}
        I'll test a single class next in examples,\\
        but you now know what a \textbf{unit} is!
    \end{tabular}
}
%~~~~~~~~~~~~~~~~~~~~~~~~~~~~~~~~~~~~~~~~~~~~%
\begin{frame}[fragile]{White box testing}
    \vspace{-3mm}
    \begin{overlayarea}{\textwidth}{0.7\textheight}
        \begin{itemize}
            \item You allow yourself in using all code in your tests\\
                  \begin{onlyenv}<1>
                      \then even private methods and members\\
                      \then implementation details are used to test interface
                  \end{onlyenv}
            \item<2-> In some languages (e.g.\ Python) easier than in others (e.g.\ C++)
            \item<only@5-> \PS{Something to consider to start working on legacy code!}\\
            \item<only@6> \alert{What happens when you change your implementation?}\\
                          \then Compilation might break\\
                          \then If not, some tests will probably fail\\
                          \then You likely need to change your tests code\\
                          \then For large refactoring, it might be boring/frustrating\\[3mm]
                          \centerline{\includegraphics[width=18mm]{Thinking}}
        \end{itemize}
        \begin{columns}[T]
            \begin{column}{0.45\textwidth}
                \begin{varblock}{example}[\textwidth]{Consider this class}<only@2-3>
                    \begin{lstlisting}[style=MyCpp]
                        // File Cup.h

                        class Cup {
                          public:
                            Cup();
                            bool IsEmpty();
                            bool Drink();
                            bool Fill();
                          private:
                            bool isEmpty_ = true;
                        };
                    \end{lstlisting}
                \end{varblock}
                \begin{varblock}{example}[\textwidth]{Consider this class}<only@4>
                    \begin{lstlisting}[style=MyCpp]
                        // File Cup.h

                        class Cup {
                          public:
                            Cup();
                            bool IsEmpty();
                            bool Drink();
                            bool Fill();
                          private:
                            friend struct CupTester;
                            bool isEmpty_ = true;
                        };
                    \end{lstlisting}
                \end{varblock}
            \end{column}
            \begin{column}{0.48\textwidth}
                \begin{varblock}{example}[\textwidth]{Test file}<only@2>
                    \begin{lstlisting}[style=MyCpp]
                        #include |+"vir/test.h"+|
                        #include |+"Cup.h"+|


                        TEST(|=Cup::Cup=|) {
                            Cup cup{};
                            VERIFY(cup.isEmpty_);
                        }
                        TEST(|=Cup::IsEmpty=|) { /*...*/ }
                        TEST(|=Cup::Fill=|)    { /*...*/ }
                        TEST(|=Cup::Drink=|)   { /*...*/ }
                    \end{lstlisting}
                \end{varblock}
                \begin{varblock}{example}[\textwidth]{Test file}<only@3>
                    \begin{lstlisting}[style=MyCpp]
                        #include |+"vir/test.h"+|
                        #include |+"Cup.h"+|
                        #define private public // UB!!!

                        TEST(|=Cup::Cup=|) {
                            Cup cup{};
                            VERIFY(cup.isEmpty_);
                        }
                        TEST(|=Cup::IsEmpty=|) { /*...*/ }
                        TEST(|=Cup::Fill=|)    { /*...*/ }
                        TEST(|=Cup::Drink=|)   { /*...*/ }
                    \end{lstlisting}
                \end{varblock}
                \begin{varblock}{example}[\textwidth]{Test file}<only@4>
                    \begin{lstlisting}[style=MyCpp]
                        #include |+"vir/test.h"+|
                        #include |+"Cup.h"+|
                        struct CupTester { /*...*/ };

                        TEST(|=Cup::Cup=|) {
                            Cup cup{};
                            CupTester tester{cup};
                            VERIFY(tester.isEmpty_);
                        }
                        TEST(|=Cup::IsEmpty=|) { /*...*/ }
                        TEST(|=Cup::Fill=|)    { /*...*/ }
                        TEST(|=Cup::Drink=|)   { /*...*/ }
                    \end{lstlisting}
                \end{varblock}
                \centering\vspace{-3mm}
                \only<2>{\alert{\scriptsize This does not compile!}}
                \only<3>{\alert{\scriptsize This relies on undefined behaviour!!!}}
            \end{column}
        \end{columns}
        \FrameRemark<2-4>{I took this example from Dave Steffen's talk at CppCon2020.}
    \end{overlayarea}
\end{frame}
%~~~~~~~~~~~~~~~~~~~~~~~~~~~~~~~~~~~~~~~~~~~~%
\begin{frame}[fragile]{Black-box testing}
    \vspace{-3mm}
    \begin{overlayarea}{\textwidth}{0.7\textheight}
        \begin{itemize}
            \item What about just testing the public interface?
        \end{itemize}
        \vspace{-3mm}
        \begin{columns}[T]
            \begin{column}{0.45\textwidth}
                \begin{varblock}{example}[\textwidth]{Consider this class}<only@1-3>
                    \begin{lstlisting}[style=MyCpp]
                        // File Cup.h

                        class Cup {
                          public:
                            Cup();
                            bool IsEmpty();
                            bool Drink();
                            bool Fill();
                          private:
                            bool isEmpty_ = true;
                        };
                    \end{lstlisting}
                \end{varblock}
                \centering%\vspace{-3mm}
                \only<1-3>{\scriptsize Source code same as before}
            \end{column}
            \begin{column}{0.48\textwidth}
                \begin{varblock}{example}[\textwidth]{Test file}<only@1>
                    \begin{lstlisting}[style=MyCpp]
                        #include |+"vir/test.h"+|
                        #include |+"Cup.h"+|

                        TEST(|=Cup::Cup=|)     { /*...*/ }
                        TEST(|=Cup::IsEmpty=|) { /*...*/ }
                        TEST(|=Cup::Fill=|)    { /*...*/ }
                        TEST(|=Cup::Drink=|)   { /*...*/ }
                    \end{lstlisting}
                \end{varblock}
                \begin{varblock}{example}[\textwidth]{Test file}<only@2>
                    \begin{lstlisting}[style=MyCpp]
                        #include |+"vir/test.h"+|
                        #include |+"Cup.h"+|

                        TEST(|=Cup::Cup=|) {
                            Cup cup{};
                            VERIFY(cup.IsEmpty());
                        }
                        TEST(|=Cup::IsEmpty=|) { /*...*/ }
                        TEST(|=Cup::Fill=|)    { /*...*/ }
                        TEST(|=Cup::Drink=|)   { /*...*/ }
                    \end{lstlisting}
                \end{varblock}
                \begin{varblock}{example}[\textwidth]{Test file}<only@3>
                    \begin{lstlisting}[style=MyCpp]
                        #include |+"vir/test.h"+|
                        #include |+"Cup.h"+|

                        TEST(|=Cup::Cup=|) {
                            Cup cup{};
                            VERIFY(cup.IsEmpty());
                        }
                        TEST(|=Cup::IsEmpty=|){
                            Cup cup{};
                            VERIFY(cup.IsEmpty());
                        }
                        TEST(|=Cup::Fill=|)    { /*...*/ }
                        TEST(|=Cup::Drink=|)   { /*...*/ }
                    \end{lstlisting}
                \end{varblock}
            \end{column}
        \end{columns}
        \begin{onlyenv}<3>
            \begin{center}
                \alert{Wait, we wrote the same code twice!}
            \end{center}
        \end{onlyenv}
        \vspace{-3mm}
        \begin{varblock}{quote}[0.8\textwidth]{The black-box conundrum}[Dave Steffen]<only@4-5>
            Fundamentally if we only test via the public interface, we have a circular-logic problem: you have to use the interface to test the interface, so at some point you have to trust something which you haven't tested.
        \end{varblock}
        \begin{varblock}{alert}[0.95\textwidth]{So, what to do then?}<only@5>
            \begin{enumerate}
                \item Ignore the problem \Remark{Be aware: Tests are not independent, one bug triggers many failures!}
                \item Declare one member function ``correct by inspection'' \Remark{and start from there}
                \item Abandon black-box testing as impractical
            \end{enumerate}
        \end{varblock}
    \end{overlayarea}
\end{frame}
%===================================================================================================%


%===================================================================================================%
\section{Test the behaviour, not the implementation}
%~~~~~~~~~~~~~~~~~~~~~~~~~~~~~~~~~~~~~~~~~~~~%
\begin{frame}[fragile]{Do not test methods! Test behaviour!}
    \vspace{-11mm}
    \begin{overlayarea}{\textwidth}{0.7\textheight}
        \begin{varblock}{example}[0.75\textwidth]{Back to our cup example:}<only@1>
            \begin{lstlisting}[style=MyCpp]
                #include |+"vir/test.h"+|
                #include |+"Cup.h"+|

                TEST(|=A new cup is empty=|)            { /*...*/ }
                TEST(|=An empty cup can be filled=|)    { /*...*/ }
                TEST(|=A filled cup is full=|)          { /*...*/ }
                TEST(|=Drinking empties a filled cup=|) { /*...*/ }
            \end{lstlisting}
        \end{varblock}
        \begin{varblock}{example}[0.75\textwidth]{Back to our cup example:}<only@2-3>
            \begin{lstlisting}[style=MyCpp]
                #include |+"vir/test.h"+|
                #include |+"Cup.h"+|

                TEST(|=A new cup is empty=|){
                    Cup cup{};
                    VERIFY(cup.IsEmpty());
                }
                TEST(|=An empty cup can be filled=|)    { /*...*/ }
                TEST(|=A filled cup is full=|)          { /*...*/ }
                TEST(|=Drinking empties a filled cup=|) { /*...*/ }
            \end{lstlisting}
        \end{varblock}
        \only<1-3>{\vspace{3mm}}
        \begin{itemize}
            \item<only@1> Test names describe the unit specifications
            \item<only@1> Suddenly, we are completely independent from implementation details
            \item<only@1> OK, so how does test code look like?
            \item<only@3> \alert{Wait!! This is exactly the same code that wasn't OK few minutes ago!?}\\
                     \then How is it that changing the name makes it fine?\\
                     \then How should this solve the black-box conundrum?
        \end{itemize}
        \begin{onlyenv}<4->
            \begin{varblock}{example}[0.6\textwidth]{The bottom line}
                We are now testing consistency of interface, not correctness of implementation!
            \end{varblock}
            \begin{varblock}{alert}[0.6\textwidth]{Let's say it once more:}<5->
                \textbf{We are now testing consistency of interface, not correctness of implementation!}
            \end{varblock}
            \begin{varblock}{}[0.9\textwidth]{It is not a philosophical statement!}<6>
                \PB{A bug that cannot be observed under any conditions is not a bug!}
            \end{varblock}
        \end{onlyenv}
    \end{overlayarea}
\end{frame}
%~~~~~~~~~~~~~~~~~~~~~~~~~~~~~~~~~~~~~~~~~~~~%
\begin{frame}[fragile]{\uncover<2>{As testing behaviour passes, the code is correct!}}
    \begin{varblock}{example}[0.9\textwidth]{Is this code correct or wrong?}
        \begin{lstlisting}[style=MyCpp]
            // File Cup.h

            class Cup {
              public:
                Cup()          { isEmpty_ = false; }
                bool IsEmpty() { return !isEmpty_; }
                bool Drink()   { /*...*/ }
                bool Fill()    { /*...*/ }
              private:
                bool isEmpty_ = false;
            };
        \end{lstlisting}
    \end{varblock}
    \vspace{-2mm}
    \begin{varblock}{example}[0.9\textwidth]{}<2>
        \begin{lstlisting}[style=MyCpp]
            #include |+"vir/test.h"+|
            #include |+"Cup.h"+|
            TEST(|=A new cup is empty=|)            { /*...*/ }
            TEST(|=An empty cup can be filled=|)    { /*...*/ }
            TEST(|=A filled cup is full=|)          { /*...*/ }
            TEST(|=Drinking empties a filled cup=|) { /*...*/ }
        \end{lstlisting}
    \end{varblock}
\end{frame}
%===================================================================================================%



%===================================================================================================%
\section{Test driven development as discipline}
%~~~~~~~~~~~~~~~~~~~~~~~~~~~~~~~~~~~~~~~~~~~~%
\AlertFrame{It is a discipline}
%~~~~~~~~~~~~~~~~~~~~~~~~~~~~~~~~~~~~~~~~~~~~%
\begin{frame}{It is a discipline}
    \begin{itemize}
        \item \alert{Do not bring it immediately into your daily work}
        \item Learn it in small exercise projects:
              \begin{itemize}
                  \item Implement a \;\PS{\texttt{Stack}}\; object
                  \item Implement a file parser
                  \item Write a 1D integration tool with few different algorithms
                  \item Simulate the 2D Ising model with a Monte Carlo algorithm
                  \item Implement a simple game (e.g.\ hangman, minesweeper)
                  \item{} [...]
              \end{itemize}
        \item \PB{Convince yourself about it}
        \item Get a feeling and choose ``your way''
    \end{itemize}
    \begin{varblock}{example}[0.78\textwidth]{Finally}
        \bfseries Get rid of the ``there-is-no-time-to-write-tests'' attitude!\\
    \end{varblock}
\end{frame}
%~~~~~~~~~~~~~~~~~~~~~~~~~~~~~~~~~~~~~~~~~~~~%
\begin{frame}{The purist original approach}
    \begin{varblock}{quote}[\textwidth]{Test driven development main rules}[Robert C. Martin]
        \begin{enumerate}[<+(1)->]
            \item You're not allowed to write any production code until you have first written a test that fails because the production code doesn't exist.
            \item You're not allowed to write more of a test than is sufficient to fail.
                  And not compiling is failing.
            \item You're not allowed to write any more production code than is sufficient to pass the currently failing test.
        \end{enumerate}
    \end{varblock}
    \begin{varblock}{example}[0.6\textwidth]{}<5>
        How would it be a world in which everything always worked one minute ago?
    \end{varblock}
\end{frame}
%~~~~~~~~~~~~~~~~~~~~~~~~~~~~~~~~~~~~~~~~~~~~%
\begin{frame}[fragile]{Writing the factorial in TDD}
    \vspace{-2mm}
    \begin{overlayarea}{\textwidth}{0.82\textheight}
        \begin{columns}[T]
            \begin{column}{0.49\textwidth}
                \begin{varblock}{example}[\textwidth]{Source file}<only@1-2>
                    \begin{lstlisting}[style=MyCpp]
                        // Empty file
                    \end{lstlisting}
                \end{varblock}
                \begin{varblock}{example}[\textwidth]{Source file}<only@3>
                    \begin{lstlisting}[style=MyCpp, belowskip=-9mm]
                        void factorial(int n){}
                    \end{lstlisting}
                    {\scriptsize\alert{Notice the function signature!}}
                \end{varblock}
                \begin{varblock}{example}[\textwidth]{Source file}<only@4-5>
                    \begin{lstlisting}[style=MyCpp, belowskip=-7mm]
                        void factorial(int n)
                        {
                          throw std::logical_error();
                        }
                    \end{lstlisting}
                    \only<4>{\scriptsize\alert{Still no need to return anything!}}
                \end{varblock}
                \begin{varblock}{example}[\textwidth]{Source file}<only@6-7>
                    \begin{lstlisting}[style=MyCpp]
                        int factorial(int n)
                        {
                          if(n<0)
                            throw std::logical_error();
                          else
                            return 1;
                        }
                    \end{lstlisting}
                \end{varblock}
                \begin{varblock}{example}[\textwidth]{Source file}<only@8-9>
                    \begin{lstlisting}[style=MyCpp]
                        int factorial(int n)
                        {
                          if(n<0)
                            throw std::logical_error();
                          else if(n==1)
                            return 1;
                          else
                            return n * factorial(n-1);
                        }
                    \end{lstlisting}
                \end{varblock}
                \begin{varblock}{example}[\textwidth]{Source file}<only@10-12>
                    \begin{lstlisting}[style=MyCpp]
                        int factorial(int n)
                        {
                          if(n<0)
                            throw std::logical_error();
                          else if(n<2)
                            return 1;
                          else
                            return n * factorial(n-1);
                        }
                    \end{lstlisting}
                \end{varblock}
                \begin{varblock}{example}[\textwidth]{Source file}<only@13>
                    \begin{lstlisting}[style=MyCpp]
                        #include |+<limits>+|

                        int factorial(int n)
                        {
                          if(n<0)
                            throw std::logical_error();
                          else if(n<2)
                            return 1;
                          else if(n>=13)
                            return INT_MAX;
                          else
                            return n * factorial(n-1);
                        }
                    \end{lstlisting}
                \end{varblock}
            \end{column}
            \hspace*{-4mm}
            \begin{column}{0.48\textwidth}
                \begin{varblock}{example}[\textwidth]{Test file}<only@1>
                    \begin{lstlisting}[style=MyCpp, belowskip=-9mm]
                        #include |+<vir/test.h>+|
                        #include |+"factorial.h"+|

                         TEST(Nothing){}
                    \end{lstlisting}
                    {\scriptsize\PS{It passes}}
                \end{varblock}
                \begin{varblock}{example}[\textwidth]{Test file}<only@2-4>
                    \begin{lstlisting}[style=MyCpp, belowskip=-10mm]
                        #include |+<vir/test.h>+|
                        #include |+"factorial.h"+|

                        TEST_CATCH(Negative_factorial,
                                   std::logical_error)
                        {
                          factorial(-17);
                        }
                    \end{lstlisting}
                    {
                        \scriptsize
                        \only<2>{\alert{It doesn't compile}}
                        \only<3>{\alert{It fails}}
                        \only<4>{\PS{It passes}}
                    }
                \end{varblock}
                \begin{varblock}{example}[\textwidth]{Test file}<only@5-6>
                    \begin{lstlisting}[style=MyCpp, belowskip=-10mm]
                        #include |+<vir/test.h>+|
                        #include |+"factorial.h"+|

                        TEST_CATCH(Negative_factorial,
                                   std::logical_error)
                        {
                          factorial(-17);
                        }
                        TEST(FactorialTest) {
                          COMPARE(1, factorial(1));
                        }
                    \end{lstlisting}
                    {\scriptsize\alt<5>{\alert{It fails}}{\PS{It passes}}}
                \end{varblock}
                \begin{varblock}{example}[\textwidth]{Test file}<only@7-8>
                    \begin{lstlisting}[style=MyCpp, belowskip=-10mm]
                        #include |+<vir/test.h>+|
                        #include |+"factorial.h"+|

                        TEST_CATCH(Negative_factorial,
                                   std::logical_error)
                        {
                          factorial(-17);
                        }
                        TEST(FactorialTest) {
                          COMPARE(1, factorial(1));
                          COMPARE(120, factorial(5));
                        }
                    \end{lstlisting}
                    {\scriptsize\alt<7>{\alert{It fails}}{\PS{It passes}}}
                \end{varblock}
                \begin{varblock}{example}[\textwidth]{Test file}<only@9-10>
                    \begin{lstlisting}[style=MyCpp, belowskip=-10mm]
                        #include |+<vir/test.h>+|
                        #include |+"factorial.h"+|

                        TEST_CATCH(Negative_factorial,
                                   std::logical_error)
                        {
                          factorial(-17);
                        }
                        TEST(FactorialTest) {
                          COMPARE(1, factorial(1));
                          COMPARE(120, factorial(5));
                          COMPARE(1, factorial(0));
                        }
                    \end{lstlisting}
                    {\scriptsize\alt<9>{\alert{It fails}}{\PS{It passes}}}
                \end{varblock}
                \begin{varblock}{example}[\textwidth]{Test file}<only@11>
                    \begin{lstlisting}[style=MyCpp, belowskip=-10mm]
                        #include |+<vir/test.h>+|
                        #include |+"factorial.h"+|

                        TEST_CATCH(Negative_factorial,
                                   std::logical_error)
                        {
                          factorial(-17);
                        }
                        TEST(FactorialTest) {
                          COMPARE(1, factorial(1));
                          COMPARE(120, factorial(5));
                          COMPARE(1, factorial(0));
                          COMPARE(479001600,
                                  factorial(12));
                        }
                    \end{lstlisting}
                    {\scriptsize\PS{It still passes}}
                \end{varblock}
                \begin{varblock}{example}[\textwidth]{Test file}<only@12>
                    \begin{lstlisting}[style=MyCpp, belowskip=-10mm]
                        #include |+<vir/test.h>+|
                        #include |+"factorial.h"+|
                        #include |+<limits>+|

                        TEST_CATCH(Negative_factorial,
                                   std::logical_error)
                        {
                            factorial(-17);
                        }
                        TEST(FactorialTest) {
                          COMPARE(1, factorial(1));
                          COMPARE(120, factorial(5));
                          COMPARE(1, factorial(0));
                          COMPARE(479001600,
                                  factorial(12));
                          COMPARE(INT_MAX,
                                  factorial(13));
                        }
                    \end{lstlisting}
                    {\scriptsize\alert{It fails}}
                \end{varblock}
                \begin{varblock}{example}[\textwidth]{Test file}<only@13>
                    \begin{lstlisting}[style=MyCpp, belowskip=-10mm]
                        #include |+<vir/test.h>+|
                        #include |+"factorial.h"+|

                        TEST_CATCH(Negative_factorial,
                                   std::logical_error)
                        {
                            factorial(-17);
                        }
                        TEST(FactorialTest) {
                          COMPARE(1, factorial(1));
                          COMPARE(120, factorial(5));
                          COMPARE(1, factorial(0));
                          COMPARE(479001600,
                                  factorial(12));
                          COMPARE(INT_MAX,
                                  factorial(13));
                        }
                    \end{lstlisting}
                    {\scriptsize\PS{It passes}}
                \end{varblock}
            \end{column}
        \end{columns}
        \FrameRemark<12->{In the code \;\CPP|INT_MAX|\; has been used for space reasons. The C++11 limit \;\CPP|std::numeric_limits<int>::max()|\; should be preferred, though.}
    \end{overlayarea}
\end{frame}
%~~~~~~~~~~~~~~~~~~~~~~~~~~~~~~~~~~~~~~~~~~~~%
\begin{frame}{How strict should you be with yourself?}
    \begin{center}
        \begin{tikzpicture}
            \node[label={[anchor=north, font=\Large, label distance=-5mm]90:{Find your way!}}] (fig)
                {\includegraphics[width=0.7\textwidth]{Find_your_way}};
        \end{tikzpicture}
        \par \alert{The more edits you do at once, the more dangerous it is!}
    \end{center}
\end{frame}
%~~~~~~~~~~~~~~~~~~~~~~~~~~~~~~~~~~~~~~~~~~~~%
\begin{frame}{Disciplines feel slow, they are not slow}
    \begin{center}
        \includegraphics[height=0.8\textheight]{TDD_not_slow}
    \end{center}
\end{frame}
%~~~~~~~~~~~~~~~~~~~~~~~~~~~~~~~~~~~~~~~~~~~~%
\begin{frame}{The Red -- Green -- Refactor approach}
    \begin{onlyenv}<1>
        \begin{varblock}{quote}[0.85\textwidth]{What Kent Beck proposed}[Ian Cooper]
            \normalsize
            \begin{description}
                \item[\alert{\textbf{Red:}}] Write a little test that doesn't work, and perhaps doesn't even compile at first.
                \item[\PS{\textbf{Green:}}] Make the test work quickly, committing whatever sins necessary in the process.
                \item[\PB{\textbf{Refactor:}}] Eliminate all of the duplication created in merely getting the test to work.
            \end{description}
        \end{varblock}
        \vspace{-2mm}
        \begin{varblock}{quote}[0.92\textwidth]{About testing the implementation}[Ian Cooper]
            Sometimes it is useful to probe the implementation details because they are difficult to understand.
            You can do that, but \textbf{delete the tests afterwards!}\\[-3mm]~
        \end{varblock}
    \end{onlyenv}
    \begin{onlyenv}<2>
        \begin{center}
            \includegraphics[width=0.9\textwidth]{TDD_Cycle}
        \end{center}
    \end{onlyenv}
\end{frame}
%~~~~~~~~~~~~~~~~~~~~~~~~~~~~~~~~~~~~~~~~~~~~%
\begin{frame}{Behaviour driven development in its original idea}
    \begin{overlayarea}{\textwidth}{0.82\textheight}
        \begin{itemize}
            \item Dan North's \URL[PB]{https://dannorth.net/introducing-bdd/}{original idea} stemmed from TDD
                  \begin{itemize}[<only@1>]
                    \item Test method names should be sentences
                    \item An expressive test name is helpful when a test fails
                    \item Determine the next most important behaviour
                    \item Requirements are behaviour, too
                    \item BDD provides a ``ubiquitous language'' for analysis
                  \end{itemize}
            \item<2-> Use cases become stories and each test implements an acceptance criterion for the story
        \end{itemize}
        \vspace{-2mm}
        \setlength{\labelsep}{\fontdimen2\font} % https://tex.stackexchange.com/a/266760/128737
        \begin{varblock}{example}[0.68\textwidth]{Example of use case as a story}<only@2>
            \textbf{Title:} Customer withdraws cash\\
            \begin{description}[~~so that]
                \setlength{\itemsep}{0pt}
                \item[\alert{As a}] customer,\\
                \item[\alert{I want}] to withdraw cash from an ATM,\\
                \item[\alert{so that}] I don't have to wait in line at the bank.
            \end{description}
        \end{varblock}
        \begin{varblock}{example}[0.8\textwidth]{Example of story acceptance criterion}<only@3>
            \textbf{Scenario:} Account is overdrawn past the overdraft limit\\
            \begin{description}
                \setlength{\itemsep}{0pt}
                \item[\alert{Given}] the account is overdrawn\\
                \item[\alert{And}] the card is valid\\
                \item[\alert{When}] the customer requests cash\\
                \item[\alert{Then}] ensure a rejection message is displayed\\
                \item[\alert{And}] ensure cash is not dispensed\\
                \item[\alert{And}] ensure the card is returned
            \end{description}
        \end{varblock}
    \end{overlayarea}
\end{frame}
%~~~~~~~~~~~~~~~~~~~~~~~~~~~~~~~~~~~~~~~~~~~~%
\begin{frame}{TDD and BDD}
    \begin{center}
        \includegraphics[width=0.62\textwidth]{Different_POV}\\[3mm]
        After all they are disciplines to deal with the same idea!
    \end{center}
\end{frame}
%===================================================================================================%


%===================================================================================================%
\section{Take-home message}
%~~~~~~~~~~~~~~~~~~~~~~~~~~~~~~~~~~~~~~~~~~~~%
\begin{frame}{Wrapping up the advice}
    \vspace{-3mm}
    \begin{enumerate}[<+->]
        \item \textcolor{bg!80!normal text.fg}{«I do not have time to write tests»} \textbf{does not make any sense!}
        \item \PP{Writing tests is an inseparable part of the code development}
        \item \alert{Tests are the only way to change code still ensuring its correctness}
        \item \PS{Test the behaviour, not the implementation}
        \item Depending on the size of the project, you can go one or another way
        \item If you do not want to strictly do TDD, still envision your tests and API usage before implementing any code
        \item Now you are aware of how writing code should be done!
        \item Testing (and hence writing) code is doing science
    \end{enumerate}
\end{frame}
%~~~~~~~~~~~~~~~~~~~~~~~~~~~~~~~~~~~~~~~~~~~~%
\begin{frame}{And now?}
    \begin{onlyenv}<1>
        \begin{center}
            \includegraphics[width=0.68\textwidth]{Wow}\\[3mm]
            If this is not your reaction right now, it will be soon!
        \end{center}
    \end{onlyenv}
    \begin{onlyenv}<2>
        \begin{center}
            \includegraphics[width=0.479\textwidth]{Old_new_way}
            \includegraphics[width=0.48\textwidth]{Time_for_change}\\[5mm]
            \Large\PS{Thanks for listening!}\\[2mm]
            \PB{Questions? Feedback?}
        \end{center}
    \end{onlyenv}
\end{frame}
%===================================================================================================%

\end{document}
